\section{Conclusiones}
El desarrollo integral de esta práctica permitió validar, tanto de forma analítica como mediante herramientas de simulación, la potencia y versatilidad de las técnicas de escalamiento en el diseño de redes eléctricas. Tras completar las fases de experimentación y análisis matemático, se desprenden las siguientes conclusiones:

\begin{itemize}
	\item \textbf{Eficacia del Escalamiento de impedancia ($k_m$):} Se comprobó que esta técnica es una herramienta indispensable para la optimización de prototipos teóricos hacia implementaciones físicas funcionales. La aplicación de un factor $k_m = 5$ demostró que es posible elevar significativamente los niveles de impedancia del circuito ($|Z_2| = 5 \times |Z_1|$) y, consecuentemente, reducir el consumo de corriente en la misma proporción sin alterar la naturaleza del sistema. Los diagramas de Bode obtenidos confirmaron de manera irrefutable que la función de transferencia es invariante ante este tipo de escalamiento, manteniendo la respuesta en magnitud y fase totalmente superponibles a las del diseño original.
	
	\item \textbf{Flexibilidad del Escalamiento en Frecuencia ($k_f$):} La experimentación demostró que es posible desplazar la respuesta dinámica de un circuito hacia cualquier punto del espectro electromagnético de manera precisa. Al utilizar un factor $k_f = 2$, se logró trasladar el pico de resonancia del sistema de $158.5\,\text{kHz}$ a exactamente $317\,\text{kHz}$, preservando la integridad de los niveles de impedancia en el nuevo punto de operación ($|Z_2| = 20\,\text{k}\Omega$) \cite{image_ead999.jpg, image_ead95f.jpg}. Esto valida que la transformación $s \rightarrow s/k_f$ es un método robusto para adaptar filtros o circuitos resonantes a bandas de frecuencia específicas sin perder su selectividad ni su factor de calidad original \cite{Explicación y Aplicación de Escalamiento.pdf}.
	
	\item \textbf{Consistencia entre la Teoría y la Praxis:} Uno de los hallazgos más relevantes fue la perfecta concordancia entre el análisis de mallas en el dominio de Laplace y los resultados arrojados por el software de simulación. La simplificación algebraica de las funciones de transferencia escaladas ($H'(s) = H(s)$ para magnitud y $H'(s) = H(s/k_f)$ para frecuencia) proporcionó el sustento teórico necesario para interpretar los cambios observados en los componentes reactivos ($L$ y $C$) \cite{image_eae4a2.jpg, image_ecb6f2.jpg}. 
	
	\item \textbf{Aplicabilidad en la Ingeniería:} En última instancia, se concluye que el escalamiento combinado es esencial en la práctica profesional. Permite al ingeniero partir de prototipos normalizados ideales y ajustarlos a componentes comerciales disponibles y a requerimientos técnicos específicos de ancho de banda y potencia. Estas metodologías no solo simplifican el proceso de diseño, sino que aseguran que el comportamiento operativo del sistema final sea predecible y fiel a los objetivos de diseño establecidos inicialmente.
\end{itemize}